\subsection{Работа с двумерными сетками}
\label{sec:prog_grid}

\subsubsection{Интерфейс работы с сетками}
Все сетки в программе наследуются от абстрактного класса \cvar{IGrid}.
Необходимые для работы с сетками таблицы узлов, свойств и связности доступны
как виртуальные методы этого класса и объявлены в заголовочном файле \cvar{grid/i_grid.hpp}. Например
\begin{itemize}
\item \cvar{Point IGrid::point(size_t ipoint)} -- получить координату $i$-ой точки,
\item \cvar{double IGrid::face_area(size_t iface)} -- площадь $i$-ой грани,
\item \cvar{std::vector<int> IGrid::tab_cell_face(size_t icell)} -- получить список индексов граней для $i$-ой ячейки.
\end{itemize}
Грани (внутренние и граничные) пронумерованы сквозным образом.
Координаты точек всегда трёхмерные. Для двумерных и одномерных задач ``лишние'' координаты приравниваются нулю.

\subsubsection{Загрузка сеток}
Для построения структурированных двумерных сеток необходимо использовать класс \cvar{RegularGrid2}:
\begin{cppcode}
// сетка 10х10 в единичном квадрате
auto grid = std::make_shared<RegularGrid2D>(0, 0, 1, 1, 10, 10);
\end{cppcode}

Неструктурированные сетки должны быть прочитаны из файла:
\begin{cppcode}
// Читаем сетку из файла /app/test_data/pebigrid.vtk
std::string fn = test_directory_file("pebigrid.vtk");
UnstructuredGrid2D grid = UnstructuredGrid2D::vtk_read(fn);
\end{cppcode}

\subsubsection{Построение неструктурированных сеток}
Строить неструктурированные сетки следует с помощью утилиты \ename{hybmesh}.
В папке \ename{test_data} корневой директории репозитория лежат скрипты построения сеток:
\begin{itemize}
\item \ename{pebigrid.py} -- pebi--сетка,
\item \ename{tetragrid.py} -- сетка, состоящая из произвольных (скошенных) трех- и четырехугольников.
\end{itemize}
Инструкции по запуску этих скриптов смотри п. \ref{sec:hybmesh}.
Эти скрипты строят равномерную неструктурированную сетку
в единичном квадрате
и записывают её в файл vtk, который впоследствии можно загрузить
в расчётную программу.
В каждом из скриптов есть параметр \cvar{N}, означающий
примерное количество ячеек в итоговой сетке.
Меняя его значение можно строить сетки разного разрешения.

Для загрузки построенной сетки в решатель необходимо файл
с сеткой поместить в каталог \ename{test_data}
и далее загрузить её в класс \cvar{UnstructuredGrid2D}.


\subsubsection{Сохранение сеточных данных}
С помощью метода \cvar{IGrid::save_vtk} сетка может быть экспортирована в vtk формат и просмотрена в Paraview.
Дополнить файл численными данными можно с помощью методов
\cvar{VtkUtils::add_cell_data} -- для данных, заданных в центрах ячеек,
\cvar{VtkUtils::add_point_data} -- для данных, заданных в узлах сетки.

\begin{cppcode}
// Создать сеточный вектор для центров ячеек
std::vector<double> field(grid.n_cells());
for (size_t icell = 0; icell < grid.n_cells(); ++icell) {
    Point p = grid.cell_center(icell);  // центр ячейки
    field[i] = p.x + 2*p.y*p.y;    // значение поля в центре ячейки
}

// сохранить сетку в файл grid_data.vtk
grid.save_vtk("grid_data.vtk");

// добавить в файл "grid_data.vtk" созданный сеточный вектор под именем "field"
VtkUtils::add_cell_data(field, "field", "grid_data.vtk");
\end{cppcode}

В \secref{sec:paraview} приведены некоторые приёмы визуализации численного решения.
